Table~\ref{tab:perseed} gives per-seed results for all $113$ reported runs,
generated from the results registry. Each row is one run, whose configuration,
data split, checkpoints and metrics are retained and whose inclusion or exclusion
is recorded in the registry. The registry holds $115$ runs; the two that do not
appear here are the learnt-prior runs whose prior network never left the
chance-level plateau, at $0.866$ and $0.899$ error on $S_0$, excluded under the
criterion stated in Section~\ref{sec:compute}.

The \textbf{Condition} column identifies which study a row belongs to.
\emph{Headline} rows use the full training set with clean labels and are the cells
of Table~\ref{tab:main}, run at five seeds ($0$--$4$). \emph{Data~$p\%$} rows are
the small-data sweep of Section~\ref{sec:smalldata}, retrained on a uniformly
random $p\%$ of the training set drawn without reading any label, at five seeds
($0$--$4$).
\emph{Noise~$\eta\%$} rows are the controlled-difficulty study of
Section~\ref{sec:difficulty}, with a fraction $\eta$ of the labels of $S$
corrupted to a uniformly random different class, also at five seeds. Rows sharing
a dataset, model, prior, objective and condition differ only in their seed.

Column $n_{\mathrm b}$ is the bound-set size actually used in the certificate, so
it varies with both the prior family and the training-data fraction. The risk
certificate is the $0$--$1$ PAC-Bayes-kl bound on the Gibbs risk at confidence
$0.99$, per certificate and per run. ``Stoch~/ PM~/ Ens'' are the stochastic,
posterior-mean and ensemble test errors, which are descriptive deployment metrics
and carry no guarantee; for the noise rows they are measured against clean test
labels, whereas the certificate in the same row refers to the corrupted-label
distribution, so the two are not comparable within those rows.

{\scriptsize\begin{longtable}[]{@{}lllllr r r r r r@{}}
\caption{Per-seed results for all 113 reported runs, one row per run.}\label{tab:perseed}\\
\toprule\noalign{}
Dataset & Model & Prior & Obj. & Condition & Seed & $n_{\text{b}}$ & Cert & Stoch & PM & Ens \\
\midrule\noalign{}\endfirsthead
\multicolumn{11}{@{}l}{\emph{Table \thetable\ continued from previous page}}\\
\toprule\noalign{}
Dataset & Model & Prior & Obj. & Condition & Seed & $n_{\text{b}}$ & Cert & Stoch & PM & Ens \\
\midrule\noalign{}\endhead
\midrule\multicolumn{11}{r@{}}{\emph{continued on next page}}\\\endfoot
\bottomrule\noalign{}\endlastfoot
CIFAR-10 & CNN & learnt & $f_{\text{bbb}}$ & headline & 0 & 25000 & 0.9796 & 0.2445 & 0.2181 & 0.2055 \\
CIFAR-10 & CNN & learnt & $f_{\text{bbb}}$ & headline & 1 & 25000 & 0.9817 & 0.2400 & 0.2163 & 0.2030 \\
CIFAR-10 & CNN & learnt & $f_{\text{bbb}}$ & headline & 2 & 25000 & 0.9804 & 0.2382 & 0.2181 & 0.2029 \\
CIFAR-10 & CNN & learnt & $f_{\text{bbb}}$ & headline & 3 & 25000 & 0.9824 & 0.2557 & 0.2214 & 0.2094 \\
CIFAR-10 & CNN & learnt & $f_{\text{bbb}}$ & headline & 4 & 25000 & 0.9762 & 0.2274 & 0.2072 & 0.1947 \\
\addlinespace[2pt]
CIFAR-10 & CNN & learnt & $f_{\text{classic}}$ & headline & 0 & 25000 & 0.3894 & 0.2865 & 0.2610 & 0.2391 \\
CIFAR-10 & CNN & learnt & $f_{\text{classic}}$ & headline & 1 & 25000 & 0.3996 & 0.2940 & 0.2692 & 0.2510 \\
CIFAR-10 & CNN & learnt & $f_{\text{classic}}$ & headline & 2 & 25000 & 0.3950 & 0.2918 & 0.2558 & 0.2408 \\
CIFAR-10 & CNN & learnt & $f_{\text{classic}}$ & headline & 3 & 25000 & 0.4015 & 0.3014 & 0.2691 & 0.2528 \\
CIFAR-10 & CNN & learnt & $f_{\text{classic}}$ & headline & 4 & 25000 & 0.3808 & 0.2771 & 0.2523 & 0.2355 \\
\addlinespace[2pt]
CIFAR-10 & CNN & learnt & $f_{\text{quad}}$ & headline & 0 & 25000 & 0.4244 & 0.2788 & 0.2564 & 0.2314 \\
CIFAR-10 & CNN & learnt & $f_{\text{quad}}$ & headline & 1 & 25000 & 0.4336 & 0.2819 & 0.2581 & 0.2403 \\
CIFAR-10 & CNN & learnt & $f_{\text{quad}}$ & headline & 2 & 25000 & 0.4282 & 0.2826 & 0.2460 & 0.2306 \\
CIFAR-10 & CNN & learnt & $f_{\text{quad}}$ & headline & 3 & 25000 & 0.4351 & 0.2907 & 0.2596 & 0.2430 \\
CIFAR-10 & CNN & learnt & $f_{\text{quad}}$ & headline & 4 & 25000 & 0.4170 & 0.2671 & 0.2425 & 0.2258 \\
\addlinespace[2pt]
CIFAR-10 & CNN & data-free & $f_{\text{quad}}$ & headline & 0 & 50000 & 0.9291 & 0.8957 & 0.9000 & 0.8942 \\
CIFAR-10 & CNN & data-free & $f_{\text{quad}}$ & headline & 1 & 50000 & 0.9302 & 0.9006 & 0.9000 & 0.8966 \\
CIFAR-10 & CNN & data-free & $f_{\text{quad}}$ & headline & 2 & 50000 & 0.9295 & 0.8960 & 0.9000 & 0.8994 \\
CIFAR-10 & CNN & data-free & $f_{\text{quad}}$ & headline & 3 & 50000 & 0.9294 & 0.8989 & 0.9000 & 0.9033 \\
CIFAR-10 & CNN & data-free & $f_{\text{quad}}$ & headline & 4 & 50000 & 0.9304 & 0.8977 & 0.9000 & 0.8984 \\
\addlinespace[2pt]
CIFAR-10 & cnn13 & learnt & $f_{\text{quad}}$ & headline & 0 & 25000 & 0.3849 & 0.2174 & 0.1794 & 0.1717 \\
CIFAR-10 & cnn13 & learnt & $f_{\text{quad}}$ & headline & 1 & 25000 & 0.3934 & 0.2247 & 0.1817 & 0.1718 \\
CIFAR-10 & cnn13 & learnt & $f_{\text{quad}}$ & headline & 2 & 25000 & 0.3849 & 0.2160 & 0.1757 & 0.1700 \\
CIFAR-10 & cnn13 & learnt & $f_{\text{quad}}$ & headline & 3 & 25000 & 0.3864 & 0.2154 & 0.1770 & 0.1682 \\
CIFAR-10 & cnn13 & learnt & $f_{\text{quad}}$ & headline & 4 & 25000 & 0.3908 & 0.2273 & 0.1818 & 0.1747 \\
\addlinespace[2pt]
CIFAR-10 & cnn15 & learnt & $f_{\text{quad}}$ & headline & 1 & 25000 & 0.5284 & 0.3380 & 0.2795 & 0.2753 \\
CIFAR-10 & cnn15 & learnt & $f_{\text{quad}}$ & headline & 2 & 25000 & 0.4425 & 0.2604 & 0.2025 & 0.2037 \\
CIFAR-10 & cnn15 & learnt & $f_{\text{quad}}$ & headline & 3 & 25000 & 0.4456 & 0.2556 & 0.1976 & 0.1955 \\
\addlinespace[2pt]
Fashion-MNIST & CNN & learnt & $f_{\text{quad}}$ & headline & 0 & 30000 & 0.1658 & 0.0924 & 0.0864 & 0.0847 \\
Fashion-MNIST & CNN & learnt & $f_{\text{quad}}$ & headline & 1 & 30000 & 0.1739 & 0.0932 & 0.0870 & 0.0838 \\
Fashion-MNIST & CNN & learnt & $f_{\text{quad}}$ & headline & 2 & 30000 & 0.1627 & 0.0908 & 0.0850 & 0.0805 \\
Fashion-MNIST & CNN & learnt & $f_{\text{quad}}$ & headline & 3 & 30000 & 0.1638 & 0.0971 & 0.0881 & 0.0837 \\
Fashion-MNIST & CNN & learnt & $f_{\text{quad}}$ & headline & 4 & 30000 & 0.1708 & 0.0920 & 0.0867 & 0.0832 \\
\addlinespace[2pt]
Fashion-MNIST & FCN & learnt & $f_{\text{bbb}}$ & headline & 0 & 30000 & 0.6616 & 0.1111 & 0.1025 & 0.1001 \\
Fashion-MNIST & FCN & learnt & $f_{\text{bbb}}$ & headline & 1 & 30000 & 0.6656 & 0.1097 & 0.1000 & 0.0996 \\
Fashion-MNIST & FCN & learnt & $f_{\text{bbb}}$ & headline & 2 & 30000 & 0.6573 & 0.1125 & 0.1018 & 0.0993 \\
Fashion-MNIST & FCN & learnt & $f_{\text{bbb}}$ & headline & 3 & 30000 & 0.6692 & 0.1139 & 0.1042 & 0.0997 \\
Fashion-MNIST & FCN & learnt & $f_{\text{bbb}}$ & headline & 4 & 30000 & 0.6622 & 0.1119 & 0.1005 & 0.0978 \\
\addlinespace[2pt]
Fashion-MNIST & FCN & learnt & $f_{\text{classic}}$ & headline & 0 & 30000 & 0.1574 & 0.1228 & 0.1148 & 0.1122 \\
Fashion-MNIST & FCN & learnt & $f_{\text{classic}}$ & headline & 1 & 30000 & 0.1593 & 0.1178 & 0.1125 & 0.1084 \\
Fashion-MNIST & FCN & learnt & $f_{\text{classic}}$ & headline & 2 & 30000 & 0.1556 & 0.1224 & 0.1166 & 0.1139 \\
Fashion-MNIST & FCN & learnt & $f_{\text{classic}}$ & headline & 3 & 30000 & 0.1586 & 0.1219 & 0.1137 & 0.1110 \\
Fashion-MNIST & FCN & learnt & $f_{\text{classic}}$ & headline & 4 & 30000 & 0.1600 & 0.1170 & 0.1103 & 0.1065 \\
\addlinespace[2pt]
Fashion-MNIST & FCN & learnt & $f_{\text{quad}}$ & headline & 0 & 30000 & 0.1740 & 0.1208 & 0.1130 & 0.1110 \\
Fashion-MNIST & FCN & learnt & $f_{\text{quad}}$ & headline & 1 & 30000 & 0.1763 & 0.1162 & 0.1089 & 0.1061 \\
Fashion-MNIST & FCN & learnt & $f_{\text{quad}}$ & headline & 2 & 30000 & 0.1713 & 0.1207 & 0.1143 & 0.1092 \\
Fashion-MNIST & FCN & learnt & $f_{\text{quad}}$ & headline & 3 & 30000 & 0.1741 & 0.1209 & 0.1124 & 0.1088 \\
Fashion-MNIST & FCN & learnt & $f_{\text{quad}}$ & headline & 4 & 30000 & 0.1776 & 0.1145 & 0.1057 & 0.1040 \\
\addlinespace[2pt]
Fashion-MNIST & FCN & data-free & $f_{\text{quad}}$ & headline & 0 & 60000 & 0.4487 & 0.2219 & 0.1979 & 0.1828 \\
Fashion-MNIST & FCN & data-free & $f_{\text{quad}}$ & headline & 1 & 60000 & 0.4463 & 0.2184 & 0.1890 & 0.1800 \\
Fashion-MNIST & FCN & data-free & $f_{\text{quad}}$ & headline & 2 & 60000 & 0.4489 & 0.2231 & 0.1898 & 0.1803 \\
Fashion-MNIST & FCN & data-free & $f_{\text{quad}}$ & headline & 3 & 60000 & 0.4478 & 0.2334 & 0.1851 & 0.1824 \\
Fashion-MNIST & FCN & data-free & $f_{\text{quad}}$ & headline & 4 & 60000 & 0.4485 & 0.2303 & 0.1836 & 0.1813 \\
\addlinespace[2pt]
MNIST & CNN & learnt & $f_{\text{quad}}$ & headline & 0 & 30000 & 0.0373 & 0.0095 & 0.0082 & 0.0083 \\
MNIST & CNN & learnt & $f_{\text{quad}}$ & headline & 1 & 30000 & 0.0369 & 0.0107 & 0.0103 & 0.0100 \\
MNIST & CNN & learnt & $f_{\text{quad}}$ & headline & 2 & 30000 & 0.0364 & 0.0124 & 0.0103 & 0.0105 \\
MNIST & CNN & learnt & $f_{\text{quad}}$ & headline & 3 & 30000 & 0.0389 & 0.0110 & 0.0096 & 0.0102 \\
MNIST & CNN & learnt & $f_{\text{quad}}$ & headline & 4 & 30000 & 0.0352 & 0.0104 & 0.0092 & 0.0087 \\
\addlinespace[2pt]
MNIST & FCN & learnt & $f_{\text{bbb}}$ & headline & 0 & 30000 & 0.2657 & 0.0191 & 0.0151 & 0.0154 \\
MNIST & FCN & learnt & $f_{\text{bbb}}$ & headline & 1 & 30000 & 0.2631 & 0.0175 & 0.0149 & 0.0148 \\
MNIST & FCN & learnt & $f_{\text{bbb}}$ & headline & 2 & 30000 & 0.2585 & 0.0186 & 0.0141 & 0.0140 \\
MNIST & FCN & learnt & $f_{\text{bbb}}$ & headline & 3 & 30000 & 0.2762 & 0.0191 & 0.0147 & 0.0145 \\
MNIST & FCN & learnt & $f_{\text{bbb}}$ & headline & 4 & 30000 & 0.2653 & 0.0192 & 0.0156 & 0.0155 \\
\addlinespace[2pt]
MNIST & FCN & learnt & $f_{\text{classic}}$ & headline & 0 & 30000 & 0.0450 & 0.0247 & 0.0206 & 0.0204 \\
MNIST & FCN & learnt & $f_{\text{classic}}$ & headline & 1 & 30000 & 0.0448 & 0.0237 & 0.0185 & 0.0193 \\
MNIST & FCN & learnt & $f_{\text{classic}}$ & headline & 2 & 30000 & 0.0450 & 0.0230 & 0.0197 & 0.0194 \\
MNIST & FCN & learnt & $f_{\text{classic}}$ & headline & 3 & 30000 & 0.0465 & 0.0235 & 0.0191 & 0.0191 \\
MNIST & FCN & learnt & $f_{\text{classic}}$ & headline & 4 & 30000 & 0.0461 & 0.0253 & 0.0182 & 0.0189 \\
\addlinespace[2pt]
MNIST & FCN & learnt & $f_{\lambda}$ & headline & 0 & 30000 & 0.0766 & 0.0214 & 0.0162 & 0.0169 \\
MNIST & FCN & learnt & $f_{\lambda}$ & headline & 1 & 30000 & 0.0767 & 0.0209 & 0.0173 & 0.0172 \\
MNIST & FCN & learnt & $f_{\lambda}$ & headline & 2 & 30000 & 0.0768 & 0.0206 & 0.0171 & 0.0163 \\
MNIST & FCN & learnt & $f_{\lambda}$ & headline & 3 & 30000 & 0.0795 & 0.0211 & 0.0167 & 0.0171 \\
MNIST & FCN & learnt & $f_{\lambda}$ & headline & 4 & 30000 & 0.0785 & 0.0214 & 0.0176 & 0.0163 \\
\addlinespace[2pt]
MNIST & FCN & learnt & $f_{\text{quad}}$ & data 10\% & 0 & 3000 & 0.1487 & 0.0623 & 0.0544 & 0.0537 \\
MNIST & FCN & learnt & $f_{\text{quad}}$ & data 10\% & 1 & 3000 & 0.1628 & 0.0630 & 0.0537 & 0.0544 \\
MNIST & FCN & learnt & $f_{\text{quad}}$ & data 10\% & 2 & 3000 & 0.1488 & 0.0628 & 0.0562 & 0.0570 \\
MNIST & FCN & learnt & $f_{\text{quad}}$ & data 10\% & 3 & 3000 & 0.1532 & 0.0631 & 0.0541 & 0.0566 \\
MNIST & FCN & learnt & $f_{\text{quad}}$ & data 10\% & 4 & 3000 & 0.1463 & 0.0648 & 0.0571 & 0.0583 \\
\addlinespace[2pt]
MNIST & FCN & learnt & $f_{\text{quad}}$ & data 20\% & 0 & 6000 & 0.1268 & 0.0452 & 0.0365 & 0.0397 \\
MNIST & FCN & learnt & $f_{\text{quad}}$ & data 20\% & 1 & 6000 & 0.1178 & 0.0482 & 0.0383 & 0.0392 \\
MNIST & FCN & learnt & $f_{\text{quad}}$ & data 20\% & 2 & 6000 & 0.1183 & 0.0480 & 0.0407 & 0.0414 \\
MNIST & FCN & learnt & $f_{\text{quad}}$ & data 20\% & 3 & 6000 & 0.1172 & 0.0466 & 0.0391 & 0.0397 \\
MNIST & FCN & learnt & $f_{\text{quad}}$ & data 20\% & 4 & 6000 & 0.1178 & 0.0459 & 0.0395 & 0.0401 \\
\addlinespace[2pt]
MNIST & FCN & learnt & $f_{\text{quad}}$ & data 50\% & 0 & 15000 & 0.0699 & 0.0297 & 0.0239 & 0.0250 \\
MNIST & FCN & learnt & $f_{\text{quad}}$ & data 50\% & 1 & 15000 & 0.0830 & 0.0315 & 0.0244 & 0.0252 \\
MNIST & FCN & learnt & $f_{\text{quad}}$ & data 50\% & 2 & 15000 & 0.0777 & 0.0280 & 0.0218 & 0.0233 \\
MNIST & FCN & learnt & $f_{\text{quad}}$ & data 50\% & 3 & 15000 & 0.0773 & 0.0321 & 0.0247 & 0.0259 \\
MNIST & FCN & learnt & $f_{\text{quad}}$ & data 50\% & 4 & 15000 & 0.0808 & 0.0305 & 0.0230 & 0.0246 \\
\addlinespace[2pt]
MNIST & FCN & learnt & $f_{\text{quad}}$ & headline & 0 & 30000 & 0.0572 & 0.0216 & 0.0170 & 0.0191 \\
MNIST & FCN & learnt & $f_{\text{quad}}$ & headline & 1 & 30000 & 0.0574 & 0.0216 & 0.0178 & 0.0177 \\
MNIST & FCN & learnt & $f_{\text{quad}}$ & headline & 2 & 30000 & 0.0572 & 0.0216 & 0.0178 & 0.0176 \\
MNIST & FCN & learnt & $f_{\text{quad}}$ & headline & 3 & 30000 & 0.0595 & 0.0218 & 0.0178 & 0.0178 \\
MNIST & FCN & learnt & $f_{\text{quad}}$ & headline & 4 & 30000 & 0.0589 & 0.0229 & 0.0176 & 0.0168 \\
\addlinespace[2pt]
MNIST & FCN & learnt & $f_{\text{quad}}$ & noise 10\% & 0 & 30000 & 0.2328 & 0.0626 & 0.0556 & 0.0393 \\
MNIST & FCN & learnt & $f_{\text{quad}}$ & noise 10\% & 1 & 30000 & 0.2335 & 0.0612 & 0.0563 & 0.0367 \\
MNIST & FCN & learnt & $f_{\text{quad}}$ & noise 10\% & 2 & 30000 & 0.2356 & 0.0589 & 0.0562 & 0.0379 \\
MNIST & FCN & learnt & $f_{\text{quad}}$ & noise 10\% & 3 & 30000 & 0.2313 & 0.0586 & 0.0527 & 0.0378 \\
MNIST & FCN & learnt & $f_{\text{quad}}$ & noise 10\% & 4 & 30000 & 0.2337 & 0.0629 & 0.0556 & 0.0386 \\
\addlinespace[2pt]
MNIST & FCN & learnt & $f_{\text{quad}}$ & noise 20\% & 0 & 30000 & 0.4246 & 0.1112 & 0.1095 & 0.0683 \\
MNIST & FCN & learnt & $f_{\text{quad}}$ & noise 20\% & 1 & 30000 & 0.4230 & 0.1196 & 0.1150 & 0.0685 \\
MNIST & FCN & learnt & $f_{\text{quad}}$ & noise 20\% & 2 & 30000 & 0.4369 & 0.1185 & 0.1150 & 0.0719 \\
MNIST & FCN & learnt & $f_{\text{quad}}$ & noise 20\% & 3 & 30000 & 0.4296 & 0.1164 & 0.1118 & 0.0681 \\
MNIST & FCN & learnt & $f_{\text{quad}}$ & noise 20\% & 4 & 30000 & 0.4305 & 0.1168 & 0.1045 & 0.0660 \\
\addlinespace[2pt]
MNIST & FCN & learnt & $f_{\text{quad}}$ & noise 40\% & 0 & 30000 & 0.6980 & 0.2649 & 0.2568 & 0.1814 \\
MNIST & FCN & learnt & $f_{\text{quad}}$ & noise 40\% & 1 & 30000 & 0.7015 & 0.2690 & 0.2582 & 0.1846 \\
MNIST & FCN & learnt & $f_{\text{quad}}$ & noise 40\% & 2 & 30000 & 0.6989 & 0.2641 & 0.2629 & 0.1813 \\
MNIST & FCN & learnt & $f_{\text{quad}}$ & noise 40\% & 3 & 30000 & 0.7048 & 0.2636 & 0.2481 & 0.1761 \\
MNIST & FCN & learnt & $f_{\text{quad}}$ & noise 40\% & 4 & 30000 & 0.7034 & 0.2619 & 0.2492 & 0.1746 \\
\addlinespace[2pt]
MNIST & FCN & data-free & $f_{\text{quad}}$ & headline & 0 & 60000 & 0.3506 & 0.0974 & 0.0562 & 0.0604 \\
MNIST & FCN & data-free & $f_{\text{quad}}$ & headline & 1 & 60000 & 0.3458 & 0.0922 & 0.0595 & 0.0569 \\
MNIST & FCN & data-free & $f_{\text{quad}}$ & headline & 2 & 60000 & 0.3496 & 0.0914 & 0.0559 & 0.0568 \\
MNIST & FCN & data-free & $f_{\text{quad}}$ & headline & 3 & 60000 & 0.3485 & 0.0884 & 0.0550 & 0.0569 \\
MNIST & FCN & data-free & $f_{\text{quad}}$ & headline & 4 & 60000 & 0.3536 & 0.0941 & 0.0582 & 0.0612 \\
\end{longtable}}

